\documentclass[11pt]{article}
\usepackage[letterpaper,margin=0.94in]{geometry}
\usepackage{amsmath,amssymb,amsthm}
\usepackage[hidelinks]{hyperref}
\newtheorem*{theorem}{Theorem}
\newcommand{\F}{\mathbb F_2}
\newcommand{\ket}[1]{\lvert #1\rangle}
\newcommand{\supp}{\operatorname{supp}}
\title{Learning the wiring of a phase-rotated cluster state}
\author{}\date{}
\begin{document}
\maketitle
\vspace{-2.8em}
Suppose a source produces copies of a linear cluster state. The qubit
labels are scrambled and each qubit has an unknown $Z$ rotation. Can we
recover the order along the chain? A two-copy measurement learns every
edge with $O(\log(n/\delta))$ copies at failure probability $\delta$.
The same protocol learns any labeled forest, even with unbounded degree.
On an arbitrary graph it returns the 2-core and all edges outside that core.

The inspiration is the path phase $\sum_i x_ix_{i+1}$ with arbitrary
one-site factors in the supplied account of an RH argument \cite{remark}.
Here a second copy of the phase cancels one-site factors at a random
set of vertices, and a finite Fourier transform reads their neighbor parities.
The algorithm uses no hypothesis about $L$-functions.

\section*{1. A two-copy measurement}
Let $G$ be an unknown simple labeled graph on $n$ vertices with adjacency
matrix $A$ over $\F$. Write $Q_G(x)=\sum_{uv\in E(G)}x_ux_v$. One input is
\[
 \ket{\psi_{G,\theta}}=2^{-n/2}\sum_{x\in\F^n}
 (-1)^{Q_G(x)} e^{i\sum_v\theta_v x_v}\ket{x}.
\]
Every one-qubit diagonal unitary has this form up to a global phase. The
basic model supplies independent measurement rounds with the same unknown
$G$. Each round supplies a pair with identical angles. The angles may
change arbitrarily between rounds. A device can prepare a copy with
$n$ Hadamards, one CZ per edge and at most $n$ one-qubit phase gates.
The learning problem gives access to the source while hiding its wiring.

At each vertex apply a CNOT from the first copy to the second and a
Hadamard to the first. Read the second result as $s_v$, the first as $t_v$.
All $n$ CNOTs run in parallel. This is an ordinary Bell measurement on
corresponding pairs of qubits.
\begin{theorem}
The string $s$ is uniform on $\F^n$. At every site with $s_v=1$,
\[
 t_v=(As)_v=\big|N(v)\cap\supp(s)\big|\pmod 2
\]
with probability one, for every choice of the angles.
\end{theorem}
\begin{proof}
Conditional on the second readout $s$, the two inputs to the phase were
$x$ and $x+s$. Their binary quadratic phases add to
$Q_G(s)+(As)\cdot x$. At site $v$ with $s_v=1$ the two unary factors
multiply to $e^{i\theta_v}$ independently of $x_v$. The remaining
wavefunction factorizes across $v$. At such a site the first qubit is
$Z^{(As)_v}\ket{+}$, so the Hadamard returns $(As)_v$. The equal
input magnitudes give $\Pr(s)=2^{-n}$.
\end{proof}
For a three-site path with middle vertex 2, $s=(1,1,0)$ returns $t_1=t_2=1$.
The third readout can be discarded. A full measurement round uses $2n$
qubits, $n$ inter-copy CNOTs, $n$ Hadamards and computational readout.

\newpage
\section*{2. A decoder that tolerates missing parities}
Take $r$ Bell rounds. For vertex $v$ store the $r$-bit fingerprint $c_v$ of
its mask bits. Store $y_v$, its $r$ output bits with entries zeroed wherever
$c_v$ is zero. The bitwise AND is denoted by $\mathbin{\&}$.
With all vertices initially active, place their labels on a queue.
\begin{enumerate}
\item Pop an active $v$. If $y_v=0$, remove it as an isolated vertex.
\item Otherwise look for an active $w\ne v$ with $y_v=c_v\mathbin{\&}c_w$.
If one exists, record $vw$, remove $v$, update
$y_w\gets y_w\mathbin{\oplus}(c_v\mathbin{\&}c_w)$, and append $w$ to the queue.
\item If neither test applies, leave $v$ active. At termination return all
recorded edges and the set of surviving vertices.
\end{enumerate}
Skipped inactive queue entries have no effect. On the success event,
the updated $y_v$ always refers to neighbors in the remaining graph.

\begin{theorem}
For $0<\delta<1$ the decoder succeeds with probability at least $1-\delta$ if
\[
 r\ \ge\ \frac{\log(2n^2/\delta)}{\log(4/3)} .
\]
In a forest it returns the exact edges. For every graph it returns the
maximal induced subgraph of minimum degree at least two as a vertex set,
and every edge with at least one endpoint outside that set. In particular
it decides whether the input graph has a cycle.
\end{theorem}
\begin{proof}
Fix the queue trajectory of a decoder that knows the graph and removes
only genuine isolated vertices and leaves. Its trajectory is independent
of the random masks. At most $n$ labels are queued initially, and at most
$n$ more are queued, one per deleted leaf. At any such test let $N$ be
the true residual neighborhood of $v$ and let $C$ be an incorrect
candidate, either empty or a singleton excluding $v$. A single round
refutes $C$ exactly when $s_v=1$ and the parity of $s$ on
$N\mathbin{\triangle}C$ is odd. Those are independent fair bits because
$N\mathbin{\triangle}C$ is nonempty and excludes $v$. Thus $C$ survives
all rounds with probability $(3/4)^r$. There are at most $n$ candidates
per test. A union bound over this fixed trajectory gives the result
without assuming independence after an adaptive deletion. When all
incorrect candidates fail, the two decoder trajectories coincide.
\end{proof}
The reference implementation uses $O(nr)$ bits for the observations and
$O(n^2)$ comparisons of $r$-bit words for the forest decoder. The latter
is a bound of $O(n^2r)$ bit operations. The randomness and all quantum
measurements precede decoding.

\paragraph{Copy lower bound.}
There are $M=n^{n-2}$ labeled trees. On $k$ copies the Hilbert space has
dimension $2^{nk}$. On uniform inputs any measurement has mean success
at most $2^{nk}/M$, since each density matrix is at most the identity and
the measurement operators sum to the identity. Even clean states therefore
require $\Omega(\log n)$ copies for constant success. Already the $n!/2$
labeled paths give the same order. Section 3 supplies the confidence
lower bound in the paired-phase model.

\newpage
\section*{3. What the measurement preserves}
Suppose each pair shares uniformly random angles, independent across sites
and freshly drawn for the next pair. Averaging a single copy gives the
maximally mixed state for every $G$. The averaged two-copy states for
\emph{all} graphs are diagonal in the same vertexwise Bell basis.
In fact, with $S=\supp(s)$ their diagonal entries are exactly
\[
 P_G(s,t)=2^{-n}2^{-(n-|S|)}\,\mathbf1_{\{t_S=(As)_S\}} .
\]
To see the diagonalization, average each phase before the Hadamards.
Only amplitudes with the same two-copy occupation at every site retain
coherence. Different $s$ have different occupation parities. At $s_v=0$
the two occupations 0 and 2 dephase, leaving a maximally mixed qubit.
At $s_v=1$ the occupation is 1 in both branches, giving the pure parity
bit of Section 1 after the Hadamard. Consequently no joint measurement
on any number of independently randomized pairs improves on Bell readout
followed by optimal classical inference.

For two specified graphs $G,H$ with equal prior probabilities let
$D=G\triangle H$ be their edgewise symmetric difference. Conditional
on $S$, the outcome laws are identical or disjoint. They are identical
exactly when every vertex of $D[S]$ has even degree. The optimal
discrimination error for $r$ pairs is therefore
\[
 \frac12\left(2^{-n}\#\{S\subseteq V: D[S]\text{ has all degrees even}\}\right)^r.
\]
If $D$ is a forest, these $S$ are exactly its independent sets. For two
graphs differing in one edge, the error is $(3/4)^r/2$. This proves
$\Omega(\log(1/\delta))$ pairs for worst-case error $\delta$ in the model
that allows phases to vary between pairs. With the counting bound, the
forest copy complexity is $\Theta(\log n+\log(1/\delta))$ for
$0<\delta\le1/3$ and $n\ge2$.

\paragraph{Drift between the two copies.}
Let the two phase vectors be $\alpha,\beta$. Conditional on $s$, the
readouts are independent bits with
\[
 \Pr(t_v\ne(As)_v\mid s)
    =\sin^2\!\big((\alpha_v+(-1)^{s_v}\beta_v)/2\big).
\]
Thus on a retained site only the phase difference matters. If the phases
are chosen independently of the masks, with
$\sin^2((\alpha_v-\beta_v)/2)\le q<1/2$, count mismatches against each
candidate across all rounds, counting unretained rows as zero.
The mean per round is at most $q/2$ for the true candidate and $1/4$
for a wrong one. Accept at most $r(1+2q)/8$ mismatches. Hoeffding's
inequality and the same ideal trajectory give failure probability at most
$2n^2\exp[-r(1-2q)^2/32]$ for independent rounds. Earlier deleted
neighbors are subtracted exactly, preserving the bit error on each row.

\paragraph{Higher degeneracy.}
For a fixed $d\ge1$ test all neighbor sets of size at most $d$. Whenever a
test succeeds remove the vertex and update those neighbors. The ideal
trajectory has at most $(d+1)n$ tests, with
$\sum_{j=0}^d\binom{n-1}{j}$ candidates each (truncate at $n-1$).
The same argument learns the fringe of the $(d+1)$-core. For
$1\le d\le n/4$ this takes $O(d\log(en/d)+\log(n/\delta))$ pairs.
Naive decoding enumerates the candidates and is polynomial for fixed $d$.

\newpage
\section*{4. A coherent-query implementation and its comparison model}
In the standard induced-subgraph parity model, a classical query on $x$
returns the single bit $Q_G(x)$. Its coherent counterpart supplies the
diagonal phase $(-1)^{Q_G(x)}$. Applying it to $\ket{+}^{\otimes n}$
prepares one clean input copy in one call. Thus the paired algorithm gives
$O(\log(n/\delta))$ coherent calls. In contrast $q$ classical binary
queries have at most $2^q$ transcripts per fixed random seed. Uniform
labeled trees imply $q\ge\log_2((1-\delta)n^{n-2})$ at worst-case error
$\delta$. This is a query separation under that access model. The copy
bound assumes a physical state source and counts neither a supplied
adjacency list nor data loaded into quantum memory.

\paragraph{Faster decoding with a callable phase gate.}
Even with unknown unary phases a forward-only oracle allows complete
fingerprints. For a chosen $S$, prepare $\ket{+}$ on $S$ and $\ket0$ on
its complement. Apply the oracle, flip all bits in $S$, apply the oracle
again with the same unary angles, undo the flips, and measure $S$ in the
$X$ basis. The unary phases cancel up to a constant. Every edge in
$G[S]$ contributes
$x_ux_v+(1-x_u)(1-x_v)=1+x_u+x_v\bmod2$. Thus all active vertices
return their degrees within $S$ modulo two. Once at the start obtain
$A\mathbf1$ by taking $S=V$. For each random mask $s$, run the experiment
on $S$ and its complement. On the complement use
$As=A\mathbf1+A(\mathbf1+s)$ to recover the missing rows.
This costs $4r+2$ forward calls and gives the complete sketches
$y_v=\bigoplus_{u\in N(v)}c_u$ with independent uniform $r$-bit labels.
In the clean model the ordinary Bell measurement already yields these
complete sketches in $2r$ calls.

With complete sketches a dictionary maps each remaining label to its
vertex. Each classifier call is one zero or dictionary lookup, and each
leaf removal requires one word XOR. Duplicate or zero fingerprints abort.
A collision among labels or with zero costs at most $n(n+1)2^{-r}/2$ in
probability. Each wrong candidate on the ideal path costs $2^{-r}$.
Hence $r=\lceil\log_2(4n^2/\delta)\rceil$ suffices. The decoder uses
$O(n)$ expected dictionary operations on $r$-bit words after $O(nr)$
processing of the measurement bits. The quantum circuit uses $O(nr)$
local gates outside the oracle.

\paragraph{Checks and prior work.}
Run \texttt{python3 bell\_forest.py} with NumPy. The checker compares the
full two-copy circuit distribution with the phase-drift formula on all 64
four-vertex graphs and reconstructs six-vertex trees from statevector
samples. Larger instances use the derived mask law. It also verifies the
phase-averaged discrimination example and the forward-only echo.
Montanaro and Shao \cite{ms} used Bell sampling to obtain $As$ for clean
graph states. Becker et al.\ \cite{becker} used low-degree pruning for
classical graph reconstruction from local messages. The result here
combines phase cancellation on selected coordinates, recovery from those
partial parities, and an exact discrimination law. Historical priority
for this combination is unverified.

\begin{thebibliography}{9}
\raggedright
\bibitem{remark} Supplied \emph{Remark}, p.~3, item 3 and footnote 3,
\texttt{remark.pdf}. See also the local \texttt{tree-cancellation.tex}.
\bibitem{ms} A. Montanaro and C. Shao, \href{https://doi.org/10.4230/LIPIcs.TQC.2022.1}
{\emph{Quantum Algorithms for Learning a Hidden Graph}}, TQC 2022.
\bibitem{becker} F. Becker et al., \href{https://arxiv.org/abs/1009.4447}
{\emph{Adding a referee to an interconnection network: What can(not) be
computed in one round}}.
\end{thebibliography}
\end{document}
